\documentclass[../../../include/open-logic-section]{subfiles}

\begin{document}

\olfileid[id]{sth}{z}{union}
\olsection{Gabungan}

\olref[sth][z][sep]{prop:intersectionsexist} memberi kita irisan. Namun, jika
kita ingin gabungan sembarang selalu ada, kita perlu menetapkan aksioma lain:

\begin{axiom}[Gabungan] Untuk setiap himpunan $A$, himpunan
$\bigcup A = \Setabs{x}{(\exists b \in A) x \in b}$ ada.
\[
	\forall A \exists U \forall x(x \in U \liff (\exists b \in A)x \in b)
\]
\end{axiom}

Aksioma ini juga dibenarkan oleh konsepsi kumulatif-iteratif. Misalkan $A$
suatu himpunan, sehingga $A$ dibentuk pada suatu tahap $S$ (menurut
\stageshier). Setiap anggota $A$ dibentuk \emph{sebelum} $S$ (menurut
\stagesacc); jadi, dengan penalaran serupa, setiap anggota dari setiap anggota
$A$ dibentuk sebelum $S$. Dengan demikian, semua himpunan \emph{tersebut}
tersedia sebelum $S$ untuk dibentuk menjadi suatu himpunan pada $S$. Himpunan
itu tepat $\bigcup A$.

\end{document}
